\part{2025}
\section{
    有积极意义的新解构——“力工梭哈理论”
} {
    \setlength{\parindent}{2em}
    网络上最近流行起“力工梭哈理论”。

力工——泛指社会层次不高的/普通的男性劳动者，这可能有潜在的导向用意；

梭哈——就是梭哈的意思，可以理解为孤注一掷，投入全部；

力工梭哈理论意思就是说：一些普通男性，把结婚生子看得太重要，极度牺牲自我，把全部积累投入到婚姻中的思维和行为模式。——潜台词是这个模式是错误和自虐的。

1，这个理论精准的否定了一些思维死板传统的人群；

2，也直接否定了支撑当年房地产经济的人群；3，也是从男性方向否定目前婚姻模式；

4，这也是要重新定义和引导男性行为和目标。从这个层次上讲，是积极的。因为：

1，现在社会环境不保护传统意义上的男性婚姻立场，如果有男性持力工梭哈思维，很可能面临人生和家庭的严重打击；

2，掏空家族钱包和背上多年房贷为婚育投入房地产的人们，可以说以近乎自虐的方式促成了如今的中国经济扭曲和困境；

3，女拳导致的女性“消极觉醒”，已经使得传统的社会共识、社会契约严重破裂，男性端的被动觉醒甚至掀桌也是必然；

4-1，旧的对男性的规训和打造性稀缺的治理方式已经逐步破产，某种意义上讲这种封建主义治理思想早该被淘汰了。

4-2，以女拳为引，导致的女性“消极觉醒”和重定义女性只是在分配上的依附性诉求，但由男女对立、责任、发展和性压抑的过度压迫导致的男性的“觉醒”和重定义可能会是整个架构的颠覆。

不过这也是一种危险的解构，因为精准命中了当今中国社会的问题——

社会治理发展管理方式的落后（核心是对国民作用和地位的理解落后倒退）与国民自身利益是有结构性矛盾的。

这个理论的解构性在于：

1，进一步解构了传统婚姻模式，男性端反婚育是真的反婚育本身，女性端大多只是想提高要价；

2，解构了传统治理学说对男性的规训控制，某种意义上讲这是迟来的男性解放，因为权责严重不对等，但男性去责任化是有风险的

3，解构体制权威和认同。因为体制既不能维护旧婚育模式的平衡，也找不到新的模式建立新平衡，其实不止婚育领域…

这种理论的出现固然极大的有人为成分，可显然是有生命力的。因为说到底是群众路线、社会主义和与时俱进长期搞的不好的问题，对，恐怕还有文科领域的长期无效发展。


source: https://user.guancha.cn/wap/content?id=1523742
}